% !TeX root = ../数模通用模板.tex
\section{模型假设}
\begin{enumerate}[label=\textbf{假设\arabic*：},leftmargin=4.8em,itemsep=6pt]
\item \textbf{以单母线聚合系统描述微网。}将光伏电源、社区负载、储能设备与外网接入点视为统一的能量交换系统，以10分钟区间平均功率建立供需平衡，忽略内部线路损耗及电压、频率的暂态变化。附件时标按区间右端点处理，整点光伏预报经插值、积分后转换为相同时间尺度的平均功率。
\item \textbf{储能效率与可用容量在研究期内保持不变。}将题面90\%的充放电效率解释为往返效率，并假设两个方向对称，取$\eta_c=\eta_d=\sqrt{0.9}$。储电量以电池内部电量计，充放电量以交流侧电量计，同一时段不同时充电和放电。调度始终满足题设容量、功率和安全储电范围，不单独计入自放电、辅助耗电及容量衰减。
\item \textbf{日初、日末储电量相等的条件仅适用于问题一。}问题一按典型日循环运行，取0:00与24:00的储电量均为6000\,kWh。问题二至问题四按连续日期运行，各方案从2025年1月1日0:00的6000\,kWh出发，使用各自的预测、规划和反馈规则完成一月份预热。每槽依据实际供需执行充放电并递推内部储电量，将1月31日24:00的实际运行状态作为2月1日0:00的初始状态。此后以前一日末储电量作为次日初值，不再施加每日初末相等的条件；不同方案的二月初始储电量可因一月运行轨迹不同而不同。
\item \textbf{将抑制不必要的功率波动和充放电切换作为运行要求。}考虑到频繁启停、短时反复换向及过大的功率变化可能增加设备运行负担，并影响电池寿命、供电运行效率和长期成本，在满足负载与储能安全边界的基础上，尽量减少碎片化充放电。各问依据模型结构，采用功率变化上限、最短运行时间、模式变化预算或功率总变差与切换惩罚实现这一要求。
\item \textbf{负载不可主动削减，外网能够正常补充供电。}社区用电需求须得到满足，不考虑需求削减及停电情形，也不另设外网购电容量约束。紧急购电仅补足常规供电与储能放电后仍存在的负载缺口，不用于充电；无法利用的剩余电量允许舍弃，其来源可以是剩余光伏或已付费的计划购电。
\item \textbf{购电决策遵守信息发布顺序，区间内反馈足够及时。}在各问允许的发布时刻，只依据此前实测记录、已知参数及已发布预报确定或调整购电量，已执行时段不再修改。假设测量和控制能够在十分钟区间内及时响应，使储能动作及紧急补购随当前实际供需确定。问题四的未来电价按逐步揭示处理，规划采用当时的价格预测，结算采用实际成交时刻的电价。
\item \textbf{近期历史误差能够反映短期预测的不确定性。}在考虑日历周期与近期状态后，假设最近至多28日的已实现预测误差可用于后续校准和情景构造，并保留整日误差路径的时间关联。该假设不要求全年误差同分布，也不要求各时段误差独立或服从正态分布；历史不足时按各问的冷启动规则处理。
\end{enumerate}

